\documentclass[draft,appendixprefix=true,headings=big]{kaobook}

\usepackage{kaoconfig}
\usepackage{showframe}
\usepackage{blindtext}
\usepackage{showlabels}

% \usepackage{minted}
% \usepackage{listings}

\begin{document}

\title{Example and documentation, v0.1.1}
\author{Federico Marotta\thanks{A \LaTeX\ lover}}
% \date{2025/12/18}

\frontmatter

\maketitle
\tableofcontents
\listoffigures
\listoftables

\mainmatter

\chapter{Layout and fonts but making it veri lonc so it overflows}
\sidetoc

\url{https://verylong.urlwhich.should~probably.most.likely.longerb.evenl.breakac-across_the.line.com}

\url{https://verylong.urlwhich.should~probably.most.likely.longerrr-across_the.line.com}

Describe the choices, the use of typearea, and the configuration options.

\section{test1}

\section{test2}

\subsection{testtest}

\setchapterimage{../../testfiles/La_Gioconda_Leonardo_Da_Vinci.jpg}
\chapter{Figures and tables}

How to use figures and tables.
There is some overlap with margin stuff: sidecaption, sidefigures and sidetables belong to both.

\begin{figure}[h]
	\includegraphics[width=.4\textwidth]{../../testfiles/La_Gioconda_Leonardo_Da_Vinci.jpg}
	\captionaboveof{table}{Mona Lisa. Leonardo da Vinci, Public domain, via \href{https://commons.wikimedia.org/wiki/File:La_Gioconda,_Leonardo_Da_Vinci.jpg}{Wikimedia Commons}.}
\end{figure}

\begin{figure}
	\includegraphics[width=.4\textwidth]{../../testfiles/La_Gioconda_Leonardo_Da_Vinci.jpg}
	\captionaboveof{table}{Mona Lisa. Leonardo da Vinci, Public domain, via \href{https://commons.wikimedia.org/wiki/File:La_Gioconda,_Leonardo_Da_Vinci.jpg}{Wikimedia Commons}.}
\end{figure}


\begin{widefigure}
	\centering
	\includegraphics[width=.4\textwidth]{../../testfiles/La_Gioconda_Leonardo_Da_Vinci.jpg}
	\captionbelowof{table}{Mona Lisa. Leonardo da Vinci, Public domain, via \href{https://commons.wikimedia.org/wiki/File:La_Gioconda,_Leonardo_Da_Vinci.jpg}{Wikimedia Commons}.}
\end{widefigure}


% \begin{kaolisting}
% 	\caption{Lasting listing impression.}
% 	\begin{minted}{Python}
% import pandas as np
% import pandas as np
% import pandas as np
% import pandas as np
% 	\end{minted}
% \end{kaolisting}

\blindtext

% \begin{lstlisting}[language=Python,caption={Wa, it works also with this caption? How is this possible?},captionpos=t]
% import numpy as pd
% import numpy as pd
% import numpy as pd
% import numpy as pd
% import numpy as pd
% import numpy as pd
% import numpy as pd
% import numpy as pd
% import numpy as pd
% import numpy as pd
% import numpy as pd
% import numpy as pd
% import numpy as pd
% import numpy as pd
% \end{lstlisting}

% \begin{kaolisting}
% \caption{Wa, it works also with this caption? How is this possible?}
% \begin{minted}{Python}
% import numpy as pd
% import numpy as pd
% import numpy as pd
% import numpy as pd
% import numpy as pd
% import numpy as pd
% import numpy as pd
% import numpy as pd
% import numpy as pd
% import numpy as pd
% import numpy as pd
% import numpy as pd
% import numpy as pd
% import numpy as pd
% \end{minted}
% \end{kaolisting}

\blindtext


\def\shortpar{Just some text to fill the margin column. It should be long enough to fill about three or four lines.}

\chapter{Margin stuff}
\sidetoc

\blindtext

Now we test sidetoc in the middle\sidetoc\ of a paragraph. what is going on?

\blindtext

\blindtext

\section{sidetoc}

Hello?

\begin{itemize}
	\item sidetoc
	\item sidepar
	\item sidenote
	\item sidecite
	\item sidecaption
	\item sidefigure
	\item sidetable
\end{itemize}

Sidetoc in horizontal\sidetoc\ mode text text text text text text text text text text text text \blindtext

\section{sidepar}

All these side items.\footnote{Are you sure you want them?}\footnote{Yes, I want them all, and I want them now.}\footnote{\blindtext\par Hello, what's the indentation?}

Testing\sidepar{\shortpar} sidepar \blindtext.

\section{sidenote}

Testing\sidenote{\shortpar} sidenote \blindtext.

Testing\sidenote{\shortpar} sidenote \blindtext.

Testing\sidenote[1em]{\shortpar} sidenote with offset \blindtext.

Testing sidenote\sidenote[0sp][1]{mark}. Asare.\\
Testing sidenote.\sidenote[0sp][1]{mark} Asare.

Testing sidenote\footnote[1]{mark}. Asare.\\
Testing sidenote.\footnote[1]{mark} Asare.

\blindtext

Ciao\fixednote{Meow}.

\section{sidefigure}

\blindtext

% \begin{sidefigure}
% 	\includegraphics[width=\textwidth]{../../testfiles/La_Gioconda_Leonardo_Da_Vinci.jpg}
% 	\caption{Mona Lisa. Leonardo da Vinci, Public domain, via \href{https://commons.wikimedia.org/wiki/File:La_Gioconda,_Leonardo_Da_Vinci.jpg}{Wikimedia Commons}.}
% \end{sidefigure}

% \blindtext

% \begin{sidetable}
% 	\caption{A useless table.}
% 	\begin{tabular}{c c c}
% 		1 & 2 & 3\\
% 		4 & 5 & 6\\
% 		7 & 8 & 9\\
% 	\end{tabular}
% \end{sidetable}

% \chapterstylelines
\chapter{References}
% \sidetoc

\section{glossaries}
Glossaries, bibliography, indices.

\section{bibliography}

\chapter{Mathematics and boxes}

% Because why not.
% Math is beautiful and should be typed beautifully.
% \begin{equation}
% 	S = \int_{t_1}^{t_2} L(q, \dot{q}, t) \, dt
% \end{equation}
% Below display.
% This line is useless, I just want to see the space between equation and surrounding text.
% \begin{description}
% 	\item[normal math] $x, y, z, \alpha, \beta, \gamma$; \textit{x}, \textit{y}, \textit{z}
% 	% \item[bold math] $\mathbfit{x}, \mathbfit{y}$; \textit{\textbf{x}}, \textit{\textbf{y}}
% 	\item[text in math] $\text{just some text}$; just some text
% 	\item[mathcal] $\mathcal{D}$, $\mathcal{R}$
% 	\item[mathscr] $\mathscr{D}$, $\mathscr{R}$
% 	\item[mathbb] $\mathbb{P}$
% \end{description}

\appendix

\chapter{Test in the app}

\backmatter

% \printbibliography

\end{document}
